% Einheit 1 von 4 – Zählen und Ergebnismengen (bank/wahrscheinlichkeit/e1.jsonl, zusammenbau v0.1)
%% TODO Zweigzeile Teil 1: was hier gelernt wird (Satz in Schülersprache aus dem Einheitstitel) – Einheit 1
\einheitenkopf[e1]{Einheit 1 von 4 · Zählen und Ergebnismengen}
\zweigzeile{ab Kl. 5, je nach Buch bis Kl. 10 (am Gymnasium ab Kl. 6) · P10}
\setcounter{aufgabe}{7}

%% TODO Titel: Ich-kann-Satz fehlt in der Bank (Platzhalter: Kettenname) – Nr. 8
%% TODO Anweisung: ein Satz über der ersten Teilaufgabe fehlt in der Bank – Nr. 8
\begin{aufgabe}{Zählen}
\begin{teile}
\teil Paul hat zwei Mützen (grau, grün) und zwei Schals (gelb, weiß). Schreibe alle Paare aus Mütze und Schal auf – erst alle mit der grauen Mütze.
\end{teile}
%% TODO Darstellung für schwach fehlt in der Bank (wahrscheinlichkeit-e1-k1-s1-v1; Abschnitt „Für schwache Schüler“)
\swz{Lena hat 3 Hosen und 2 Pullover. Wie viele Kombinationen aus Hose und Pullover gibt es?}{}
%% TODO Darstellung für schwach fehlt in der Bank (wahrscheinlichkeit-e1-k1-s1-v2; Abschnitt „Für schwache Schüler“)
\swz{Beim Bäcker gibt es 2 Brotsorten und 4 Beläge. Wie viele belegte Brote aus einer Sorte und einem Belag gibt es?}{}
%% TODO Darstellung für schwach fehlt in der Bank (wahrscheinlichkeit-e1-k1-s1-v3; Abschnitt „Für schwache Schüler“)
\swz{Ein Menü besteht aus einer von 3 Vorspeisen und einem von 3 Hauptgerichten. Wie viele Menüs gibt es?}{}
%% TODO Darstellung für schwach fehlt in der Bank (wahrscheinlichkeit-e1-k1-s1-v4; Abschnitt „Für schwache Schüler“)
\swz{Vom Haus zum Park führen 2 Wege, vom Park zur Schule 3 Wege. Wie viele Wege vom Haus über den Park zur Schule gibt es?}{}
%% TODO Darstellung für schwach fehlt in der Bank (wahrscheinlichkeit-e1-k1-s1-v5; Abschnitt „Für schwache Schüler“)
\swz{Ein Fahrrad gibt es mit 4 Rahmenfarben und 2 Sattelfarben. Wie viele Farbkombinationen gibt es?}{}
\swfrage{Was bleibt gleich, was ändert sich?}
%% TODO Darstellung für schwach fehlt in der Bank (wahrscheinlichkeit-e1-k1-s2-v1; Abschnitt „Für schwache Schüler“)
\swz{Tom hat 5 Hosen und 7 T-Shirts. Wie viele Kombinationen gibt es?}{}
%% TODO Darstellung für schwach fehlt in der Bank (wahrscheinlichkeit-e1-k1-s3-v1; Abschnitt „Für schwache Schüler“)
\swz[3]{Anna, Ben und Cem stellen sich in eine Reihe. Schreibe alle Reihenfolgen auf. Wie viele sind es?}{}
%% TODO Darstellung für schwach fehlt in der Bank (wahrscheinlichkeit-e1-k1-s4-v1; Abschnitt „Für schwache Schüler“)
\swz{5 Kinder setzen sich auf 5 Stühle in einer Reihe. Wie viele Sitzordnungen gibt es?}{}
%% TODO Darstellung für schwach fehlt in der Bank (wahrscheinlichkeit-e1-k1-s5-v1; Abschnitt „Für schwache Schüler“)
\swz{Bilde aus den Ziffern 3, 9 und 5 – jede genau einmal – die größte und die kleinste dreistellige Zahl.}{}
%% TODO Darstellung für schwach fehlt in der Bank (wahrscheinlichkeit-e1-k1-s6-v1; Abschnitt „Für schwache Schüler“)
\swz{Zwei gleiche Spielsteine mit einer roten und einer blauen Seite werden gleichzeitig geworfen. Wie viele verschiedene Ergebnisse sind möglich? \hfill (P10 2017 OS)}{}
%% TODO Darstellung für schwach fehlt in der Bank (wahrscheinlichkeit-e1-k1-s7-v1; Abschnitt „Für schwache Schüler“)
\swz[3]{Ein Würfel wird zweimal geworfen. Schreibe alle Paare auf, bei denen der zweite Wurf eine 4 zeigt. \hfill (P10 2020 OS)}{}
%% TODO Darstellung für schwach fehlt in der Bank (wahrscheinlichkeit-e1-k1-s8-v1; Abschnitt „Für schwache Schüler“)
\swz{Die vier Punkte A, B, C, D liegen so, dass A, B und C auf einer Geraden liegen. Wie viele Dreiecke mit drei dieser Punkte als Ecken gibt es?}{}
\swz[3]{Ein Fahrradschloss hat einen dreistelligen Code aus den Ziffern 2, 4 und 8; jede Ziffer kommt genau einmal vor. Wie viele Codes sind möglich? \hfill (P10 2017 OS)}{}
\end{aufgabe}

%% TODO Titel: Ich-kann-Satz fehlt in der Bank (Platzhalter: Kettenname) – Nr. 9
%% TODO Anweisung: ein Satz über der ersten Teilaufgabe fehlt in der Bank – Nr. 9
\begin{aufgabe}{Ergebnisse zu „mindestens einmal …“ aufzählen}
%% TODO Darstellung für schwach fehlt in der Bank (wahrscheinlichkeit-e1-k2-s1-v1; Abschnitt „Für schwache Schüler“)
\swz[3]{Eine Münze wird zweimal geworfen. Schreibe alle Ergebnisse auf, bei denen mindestens einmal Kopf fällt.}{}
\end{aufgabe}

%% TODO Titel: Ich-kann-Satz fehlt in der Bank (Platzhalter: Kettenname) – Nr. 10
%% TODO Anweisung: ein Satz über der ersten Teilaufgabe fehlt in der Bank – Nr. 10
\begin{aufgabe}{Zählen – Fehler finden, Begründen, Darstellungswechsel, Anwendung}
\swz[3]{Ole wirft zwei Münzen und schreibt: „Mögliche Ergebnisse sind (Kopf; Kopf), (Kopf; Zahl) und (Zahl; Zahl), also 3.“ Finde den Fehler.}{}
\swfrage{Sara schreibt alle Eisbecher aus einer Sorte (Zitrone, Kirsch) und einer Waffel (klein, groß) geordnet auf: Zitrone–klein, Zitrone–groß, Kirsch–klein, Kirsch–groß. Begründe, warum keine Möglichkeit fehlt.}
\swz[3]{Der Baum zeigt die Wahl eines Getränks und eines Kuchens. Schreibe alle Möglichkeiten als Liste auf.}{\baumzwei{Saft/,Tee/}{Apfel/,Nuss/}{Apfel/,Nuss/}}
\swz[3]{Ein Kofferschloss hat drei Ringe mit je den Ziffern 0 bis 9. Du probierst jede Sekunde einen Code. Wie lange dauert es höchstens, bis du alle Codes durchprobiert hast?}{}
\end{aufgabe}

\uebersichtskasten{Zählen: Gibt es für die erste Wahl 4 Möglichkeiten und für die zweite 3, gibt es zusammen 4 · 3 = 12 Kombinationen. \\ Anordnungen: 4 verschiedene Dinge in eine Reihe bringen: 4 · 3 · 2 · 1 = 24 Reihenfolgen. \\ Aufzählen: erst alle Ergebnisse mit dem ersten Element, dann alle mit dem zweiten – so fehlt nichts. Zwei Münzen: (Z; Z), (Z; W), (W; Z), (W; W) – (Z; W) und (W; Z) sind zwei Ergebnisse. \\ Größte Zahl aus den Ziffern 4, 8, 1: Ziffern von groß nach klein → 841}
